% Theorem numbers for Efimov_Large_Categories refer to arXiv:2405.12169v4.
	\chapter{Examples of dualizable categories}
	\label[chapter]{chap:Examples}

	% Source: Talk 3 (Max Glaser, 27.10.2026)
	In this chapter we discuss several sources of dualizable (and, more generally, compactly assembled) categories:
	\begin{itemize}
		\item kernels of homological epimorphisms of ring spectra, and in particular almost ring theory (\Cref{sec:Homological_Epimorphisms,sec:Almost_Ring_Theory});
		\item continuous posets, as a class of non-stable compactly assembled categories (\Cref{sec:Continuous_Posets});
		\item rigid symmetric monoidal categories, such as sheaves on compact Hausdorff spaces and categories of nuclear modules in condensed mathematics (\Cref{sec:Rigid_Categories,sec:Nuclear_Modules}).
	\end{itemize}

	\section{Kernels of homological epimorphisms}
	\label[section]{sec:Homological_Epimorphisms}

	% Source: Clausen, Lecture 2 (18.07.2023); Mathew, Lecture 3 (MPIM Bonn, 10.09.2024)
	Let $A \to B$ be a map of $\bbE_1$-rings, and consider the base change functor $- \otimes_A B\colon \D(A) \to \D(B)$ between the categories of left modules, with right adjoint the restriction functor $\res\colon \D(B) \to \D(A)$. Recall from \Cref{def:Homological_Epimorphism} that $A \to B$ (or rather $\Perf(A) \to \Perf(B)$) is a \emph{homological epimorphism} if $- \otimes_A B$ is a localization, i.e.\ if $\res$ is fully faithful. Since $\res$ preserves colimits, this is equivalent to the counit $B \otimes_A \res(N) \to N$ being an equivalence for all $N \in \D(B)$, and since both sides preserve colimits in $N$ and $\D(B)$ is generated by $B$, this is in turn equivalent to the multiplication map
	\[
		B \otimes_A B \to B
	\]
	being an equivalence. In the language of Clausen, $A \to B$ is then called \emph{idempotent}.

	\begin{proposition}
		\label[proposition]{prop:Kernel_Homological_Epi}
		Let $A \to B$ be a homological epimorphism of $\bbE_1$-rings, and put
		\[
			\D(A \to B) := \ker\big(- \otimes_A B\colon \D(A) \to \D(B)\big).
		\]
		Then $\D(A \to B)$ is dualizable, and the sequence $\D(A \to B) \hookrightarrow \D(A) \to \D(B)$ is a strongly continuous localization sequence. In particular,
		\[
			K^{\cont}(\D(A \to B)) \simeq \fib\big(K(A) \to K(B)\big).
		\]
	\end{proposition}
	\prooffor{prop:Kernel_Homological_Epi}
	\begin{proof}
		The functor $- \otimes_A B$ is a localization by assumption, and it is strongly continuous since its right adjoint $\res$ preserves colimits. Hence $\D(A \to B) \hookrightarrow \D(A) \to \D(B)$ is a strongly continuous localization sequence (\Cref{lem:Strong_Continuity_Localization_Sequence}), and $\D(A \to B)$ is dualizable by \Cref{prop:Localization_Sequence_Dualizable}. The final statement follows from the definition of $K^{\cont}$, see \Cref{sec:Continuous_K_Theory_Intro}.
	\end{proof}

	\begin{proposition}
		\label[proposition]{prop:Every_Dualizable_Kernel}
		Every dualizable category is of the form $\D(E \to E')$ for some homological epimorphism $E \to E'$ of $\bbE_1$-rings.
	\end{proposition}
	\prooffor{prop:Every_Dualizable_Kernel}
	\begin{proof}
		Let $C$ be dualizable. Choose a set $S$ of representatives of the isomorphism classes of $C^{\omega_1}$, a regular cardinal $\kappa > |S|$, and put $G := \bigoplus_{x \in S} x$. As a $\kappa$-small sum of $\kappa$-compact objects, $G$ is $\kappa$-compact. Let $A \subseteq C^{\kappa}$ be the thick subcategory generated by $G$; it contains $C^{\omega_1}$, since every object of $C^{\omega_1}$ is a retract of $G$. The functor $\hat{y}\colon C \to \Ind(C^{\kappa})$ takes values in $\Ind(C^{\omega_1}) \subseteq \Ind(A)$ (\Cref{thm:Compactly_Assembled_Omega_1}, see also \cite[Corollary~2.14]{Efimov_Large_Categories}), and hence is a fully faithful left adjoint of $\colim\colon \Ind(A) \to C$. Its essential image $\hat{y}(C)$ is closed under colimits, with right orthogonal $\ker(\colim\colon \Ind(A) \to C)$. Since the latter is closed under colimits in $\Ind(A)$, the localization
		\[
			Q\colon \Ind(A) \to \Ind(A)/\hat{y}(C) \simeq \ker(\colim)
		\]
		has a colimit-preserving right adjoint and thus preserves compact objects. As $Q$ is essentially surjective and $\Ind(A)$ is compactly generated by $G$, the category $\ker(\colim)$ is compactly generated by the single object $Q(G)$. Hence $\ker(\colim) \simeq \Ind(B)$ for $B$ the thick subcategory generated by $Q(G)$, and $Q$ is $\Ind$ of a homological epimorphism $A \to B$ with kernel $\hat{y}(C) \simeq C$; this is the Calkin sequence of \Cref{thm:Lurie_Dualizable_Characterization} with $C^{\kappa}$ replaced by $A$. Finally, by Morita theory \cite[Theorem~7.1.2.1]{lurie2016HA} we have $\Ind(A) \simeq \D(E)$ for $E := \End(G)\catop$ and $\Ind(B) \simeq \D(E')$ for $E' := \End(Q(G))\catop$, under which $Q$ corresponds to $- \otimes_E E'$, since both functors preserve colimits and agree on the generator. So $C \simeq \D(E \to E')$.
	\end{proof}
	% In Clausen's lectures this was attributed to Efimov; we have not found it stated in this form in Efimov's papers. The proof above was written with AI assistance and should be checked.

	There is a convenient criterion for a map $A \to B$ with fiber $I$ to be a homological epimorphism. Recall that a non-unital ring (spectrum) $I$ is called \emph{h-unital} if the bar construction gives an equivalence $I \otimes_I I \iso I$.

	\begin{theorem}[Wodzicki, Tamme]
		\label[theorem]{thm:H_Unital}
		For a non-unital ring spectrum $I$, the following are equivalent:
		\begin{enumerate}
			\item $I$ is h-unital;
			\item the map $\S \oplus I \to \S$ is a homological epimorphism;
			\item for every fiber sequence $I \to A \to B$ in which $A \to B$ is a map of unital ring spectra, the map $A \to B$ is a homological epimorphism.
		\end{enumerate}
	\end{theorem}

	Conversely, if $A \to B$ is a homological epimorphism of connective ring spectra, then $\fib(A \to B)$ is h-unital.

	\begin{example}
		\label[example]{ex:Homological_Epis}
		The following maps are homological epimorphisms, and hence give rise to dualizable categories $\D(A \to B)$:
		\begin{enumerate}
			\item A localization $A \to A[S^{-1}]$ of a commutative ring at a multiplicative subset. If $A$ is Noetherian, every homological epimorphism out of $A$ is glued from such localizations, and in this case $\D(A \to B)$ is compactly generated.
			\item A surjection $A \twoheadrightarrow B$ of perfect $\F_p$-algebras: one has $B \otimes^{\bL}_A B \simeq B$, since perfect $\F_p$-algebras have no higher $\mathrm{Tor}$-groups. Similarly, if $A$ is an integral perfectoid ring and $B$ is a perfect $\F_p$-algebra, then any surjection $A \twoheadrightarrow B$ is a homological epimorphism.
			\item (Wodzicki) If $A$ is a unital C$^*$-algebra and $I \subseteq A$ a closed two-sided $*$-ideal, then $A \to A/I$ is a homological epimorphism. For $A = B(\Hh)$ and $I = \Kk(\Hh)$, this gives the category $\D(\Kk)$ from \Cref{ex:Compact_Operators_Intro}.
			\item (Almost ring theory) If $A$ is a ring and $I \subseteq A$ is a two-sided ideal with $I^2 = I$ which is flat as a left or right $A$-module, then $A \to A/I$ is a homological epimorphism. We discuss this in detail in \Cref{sec:Almost_Ring_Theory}.
		\end{enumerate}
	\end{example}
	% TODO: references for the claims in Example ex:Homological_Epis(1)-(2) (Clausen, Lecture 2 and Mathew, Lecture 3 state them without proof).

	\section{Almost ring theory}
	\label[section]{sec:Almost_Ring_Theory}

	% Source: Talk 3 (Max Glaser, 27.10.2026)
	Almost ring theory (Faltings, Gabber--Ramero \cite{Gabber_Ramero_Almost}) provides one of the most natural sources of dualizable categories which are not compactly generated. Throughout this section, let $A$ be an (associative, unital) ring and let $I \subseteq A$ be a two-sided ideal such that
	\begin{itemize}
		\item $I^2 = I$, and
		\item $I$ is flat as a right $A$-module.
	\end{itemize}
	We work with left $A$-modules, and write $A\text{-}\Mod$ for the abelian category of left $A$-modules and $\D(A)$ for its derived category.

	The forgetful functor $(A/I)\text{-}\Mod \hookrightarrow A\text{-}\Mod$ is fully faithful, with essential image the $A$-modules $M$ with $IM = 0$. Since $I^2 = I$, this subcategory is closed under extensions: if $0 \to M' \to M \to M'' \to 0$ is exact with $IM' = 0 = IM''$, then $IM = I^2M \subseteq I M' = 0$. As it is clearly closed under subobjects and quotients, it is a \emph{Serre subcategory}.

	\begin{definition}
		\label[definition]{def:Almost_Modules}
		The category of \emph{$(A,I)$-almost modules} is the Serre quotient
		\[
			(A,I)^a\text{-}\Mod := A\text{-}\Mod/(A/I)\text{-}\Mod.
		\]
	\end{definition}

	\begin{example}
		\label[example]{ex:Almost_Modules}
		\leavevmode
		\begin{enumerate}
			\item If $I = A$, then $(A/I)\text{-}\Mod = 0$ and $(A,A)^a\text{-}\Mod = A\text{-}\Mod$.
			\item The motivating examples come from non-discrete valuation rings: let $A$ be a valuation ring of rank one whose value group is not discrete, e.g.\ $A = \Oo_{\C_p}$, and let $I = \mm$ be its maximal ideal. Then $\mm^2 = \mm$, since every element of $\mm$ is a product of two elements of $\mm$, and $\mm$ is flat since it is torsion-free over a valuation ring. An $A$-module is called \emph{almost zero} if it is killed by $\mm$; for instance, the residue field $A/\mm$ becomes zero in the almost category. This is the setting of Faltings' almost purity theorem in $p$-adic Hodge theory.
		\end{enumerate}
	\end{example}

	\subsection*{The derived picture}
	We have an adjoint triple
	\[
		\begin{tikzcd}[column sep=large]
			\D(A) \arrow[r, shift left=2.5, "L"] \arrow[r, shift right=2.5, "R"'] & \D(A/I) \arrow[l, hook', "\res" description]
		\end{tikzcd}
	\]
	where $\res$ denotes restriction of scalars, $L = A/I \otimes^{\bL}_A -$ and $R = \bR\Hom_A(A/I,-)$.

	\begin{lemma}
		\label[lemma]{lem:Almost_Restriction_FF}
		The functor $\res\colon \D(A/I) \to \D(A)$ is fully faithful. In particular, $L$ is a localization and $A \to A/I$ is a homological epimorphism.
	\end{lemma}
	\prooffor{lem:Almost_Restriction_FF}
	\begin{proof}
		As explained at the beginning of \Cref{sec:Homological_Epimorphisms}, it suffices to show that the multiplication map $A/I \otimes^{\bL}_A A/I \to A/I$ is an equivalence. Using the short exact sequence $0 \to I \to A \to A/I \to 0$ of right $A$-modules, we find
		\begin{align*}
			\mathrm{Tor}^A_0(A/I,A/I) &= A/I, \\
			\mathrm{Tor}^A_1(A/I,A/I) &\cong I/I^2 = 0, \\
			\mathrm{Tor}^A_k(A/I,A/I) &\cong \mathrm{Tor}^A_{k-1}(I,A/I) = 0 \quad \text{for } k \geq 2,
		\end{align*}
		where the last vanishing uses flatness of $I$ as a right $A$-module.
	\end{proof}

	\begin{lemma}
		\label[lemma]{lem:Almost_Essential_Image}
		The essential image of $\res\colon \D(A/I) \hookrightarrow \D(A)$ is the full subcategory
		\[
			\D_{(A/I)\text{-}\Mod}(A) := \{X \in \D(A) \mid IH_n(X) = 0 \text{ for all } n \in \Z\}
		\]
		of complexes whose homology lies in the Serre subcategory $(A/I)\text{-}\Mod$. Moreover, the kernel of $L$ consists of those $X \in \D(A)$ with $X \simeq I \otimes^{\bL}_A X$, and $I \otimes^{\bL}_A X \in \ker(L)$ for every $X \in \D(A)$.
	\end{lemma}
	\prooffor{lem:Almost_Essential_Image}
	\begin{proof}
		Since $\res$ is fully faithful with left adjoint $L$, its essential image consists of those $X$ for which the unit $X \to \res L(X)$ is an equivalence. Tensoring the short exact sequence $0 \to I \to A \to A/I \to 0$ of $A$-bimodules with $X$ gives a cofiber sequence
		\[
			I \otimes^{\bL}_A X \to X \to \res L(X),
		\]
		so $X$ lies in the essential image if and only if $I \otimes^{\bL}_A X \simeq 0$. Since $I$ is flat as a right module, $H_n(I \otimes^{\bL}_A X) \cong I \otimes_A H_n(X)$. Finally, for a left $A$-module $M$, we have $I \otimes_A M = 0$ if and only if $IM = 0$: the multiplication map $I \otimes_A M \to IM$ is surjective, and conversely if $IM = 0$ then $M$ is an $A/I$-module and $I \otimes_A M \cong (I \otimes_A A/I) \otimes_{A/I} M \cong I/I^2 \otimes_{A/I} M = 0$.

		For the second statement, $L(I \otimes^{\bL}_A X) \simeq (A/I \otimes^{\bL}_A I) \otimes^{\bL}_A X$, and $A/I \otimes^{\bL}_A I \simeq \fib(A/I \to A/I \otimes^{\bL}_A A/I) \simeq 0$ by \Cref{lem:Almost_Restriction_FF}. The characterization of $\ker(L)$ follows from the cofiber sequence above.
	\end{proof}

	\begin{proposition}
		\label[proposition]{prop:Almost_Derived_Category}
		There is an equivalence
		\[
			\D((A,I)^a\text{-}\Mod) \simeq \ker\big(L\colon \D(A) \to \D(A/I)\big),
		\]
		and $\ker(L) = \D(A \to A/I)$ is dualizable. More precisely, $\D((A,I)^a\text{-}\Mod) \hookrightarrow \D(A) \xrightarrow{L} \D(A/I)$ is a strongly continuous localization sequence, and $\D((A,I)^a\text{-}\Mod)$ is also the kernel of $L$ computed in $\Pr^{\dual}$.
	\end{proposition}
	\prooffor{prop:Almost_Derived_Category}
	\begin{proof}
		Since $(A/I)\text{-}\Mod \subseteq A\text{-}\Mod$ is a Serre subcategory of a Grothendieck abelian category which is closed under colimits, the derived category of the Serre quotient is the Verdier quotient
		\[
			\D((A,I)^a\text{-}\Mod) \simeq \D(A)/\D_{(A/I)\text{-}\Mod}(A).
		\]
		Indeed, the quotient functor $A\text{-}\Mod \to (A,I)^a\text{-}\Mod$ is exact and has a fully faithful right adjoint, and for such functors this identification is proved in \cite[Section~7.7]{Krause_Localization}; see also \cite[Example~2.26]{Efimov_Large_Categories}. It thus remains to identify the Verdier quotient $\D(A)/\D_{(A/I)\text{-}\Mod}(A)$ with $\ker(L)$. Consider the functor $q := I \otimes^{\bL}_A -\colon \D(A) \to \ker(L)$, which lands in $\ker(L)$ by \Cref{lem:Almost_Essential_Image}. For $K \in \ker(L)$ and $X \in \D(A)$, applying $\Map(K,-)$ to the cofiber sequence $I \otimes^{\bL}_A X \to X \to \res L(X)$ and using $\Map(K,\res L X) \simeq \Map(LK,LX) \simeq *$ shows that $\Map(K,qX) \iso \Map(K,X)$. So $q$ is right adjoint to the fully faithful inclusion $\ker(L) \hookrightarrow \D(A)$, and hence $q$ exhibits $\ker(L)$ as the localization of $\D(A)$ at the morphisms inverted by $q$ \cite[Proposition~5.2.7.12]{lurie2009HTT}. Since $q$ is exact, these are the morphisms whose cofiber lies in $\ker(q)$, which is $\D_{(A/I)\text{-}\Mod}(A)$ by \Cref{lem:Almost_Essential_Image}. Hence $\ker(L) \simeq \D(A)/\D_{(A/I)\text{-}\Mod}(A)$.

		By \Cref{prop:Kernel_Homological_Epi} and \Cref{lem:Almost_Restriction_FF}, $\ker(L)$ is dualizable and its inclusion into $\D(A)$ is strongly continuous; explicitly, its right adjoint is $X \mapsto I \otimes^{\bL}_A X$ by \Cref{lem:Almost_Essential_Image}. To see that $\ker(L)$ is also the kernel in $\Pr^{\dual}$, let $E$ be dualizable and $F\colon E \to \D(A)$ strongly continuous with $LF \simeq 0$. Then $F$ factors uniquely through a colimit-preserving functor $F'\colon E \to \ker(L)$, and $F'$ is strongly continuous since its right adjoint is the composite of the colimit-preserving functors $\ker(L) \hookrightarrow \D(A) \xrightarrow{F^R} E$.
	\end{proof}

	\begin{remark}
		\label[remark]{rem:Almost_General}
		In the talk, the last part of \Cref{prop:Almost_Derived_Category} was deduced from the general fact that the inclusion $\Pr^{\dual} \hookrightarrow \PrL_{\st}$ preserves pullbacks along localizations (\Cref{chap:Calkin_Limits_Colimits}). Efimov \cite[Example~2.26]{Efimov_Large_Categories} shows that $\D((A,I)^a\text{-}\Mod)$ is dualizable even without any flatness assumption on $I$: in general, $\D(A) \to \D_{(A/I)\text{-}\Mod}(A)$ is still a strongly continuous localization, identified with $- \otimes_A B$ for a connective DG ring $B$ with $H_0(B) = A/I$, and $B = A/I$ if $I$ is flat as a right or left $A$-module.
	\end{remark}

	The following is a special case of Keller's example (\Cref{ex:Keller}), and shows that almost module categories are typically very far from being compactly generated.

	\begin{proposition}
		\label[proposition]{prop:Almost_No_Compacts}
		If $I \subseteq \mathrm{Jac}(A)$ is contained in the Jacobson radical of $A$, then $\D((A,I)^a\text{-}\Mod)$ has no nonzero compact objects.
	\end{proposition}
	\prooffor{prop:Almost_No_Compacts}
	\begin{proof}
		By \Cref{prop:Almost_Derived_Category}, the inclusion $\ker(L) \hookrightarrow \D(A)$ is strongly continuous and hence preserves compact objects. So it suffices to show: if $P \in \Perf(A)$ satisfies $A/I \otimes^{\bL}_A P \simeq 0$, then $P \simeq 0$.

		Suppose $P \neq 0$ and let $n$ be minimal with $H_n(P) \neq 0$. We may represent $P$ by a bounded complex of finitely generated projective modules concentrated in homological degrees $\geq n$: starting from any bounded complex $P_\bullet$ of finitely generated projectives, if $a < n$ is the lowest degree with $P_a \neq 0$, then $P_{a+1} \to P_a$ is surjective as $H_a(P) = 0$, hence split, and we may replace $P_{a+1}$ by its kernel (again finitely generated projective) and remove $P_a$. In this representation, $H_n(P) = \coker(P_{n+1} \to P_n)$ is a finitely generated $A$-module, and by right exactness of the tensor product
		\[
			0 = H_n(A/I \otimes^{\bL}_A P) \cong A/I \otimes_A H_n(P) \cong H_n(P)/IH_n(P).
		\]
		So $H_n(P) = IH_n(P) \subseteq \mathrm{Jac}(A)H_n(P)$, and Nakayama's lemma implies $H_n(P) = 0$, a contradiction.
	\end{proof}

	\begin{example}
		\label[example]{ex:Almost_Modules_Over_Cp}
		In the situation of \Cref{ex:Almost_Modules}(2), $A = \Oo_{\C_p}$ is local with maximal ideal $\mm = \mathrm{Jac}(A)$. Hence $\D((\Oo_{\C_p},\mm)^a\text{-}\Mod)$ is a dualizable category without nonzero compact objects.
	\end{example}

	\section{Continuous posets}
	\label[section]{sec:Continuous_Posets}

	% Source: Talk 3 (Max Glaser, 27.10.2026); proofs completed with generative AI (Claude)
	We now study a class of non-stable compactly assembled categories in detail, which was already briefly discussed in \Cref{ex:Continuous_Posets}. They will reappear in \Cref{chap:Sheaves}, where we study sheaves on continuous posets. See \cite{Gierz_Continuous_Lattices} for a comprehensive treatment.

	Throughout, let $P$ be a partially ordered set, regarded as a category.

	\begin{lemma}
		\label[lemma]{lem:Ind_Of_Poset}
		The category $\Ind(P)$ is (equivalent to) a partially ordered set. Explicitly, every object of $\Ind(P)$ is of the form $\quotcolim_{s \in S} s$ for a directed subset $S \subseteq P$, and for directed subsets $S,T \subseteq P$ we have
		\begin{align*}
			\Map_{\Ind(P)}\Big(\quotcolim_{s \in S} s, \quotcolim_{t \in T} t\Big) &= \lim_{s \in S\catop}\colim_{t \in T}\Map_P(s,t) \\
			&\simeq \begin{cases} * & \text{if for all } s \in S \text{ there is } t \in T \text{ with } s \leq t, \\ \emptyset & \text{otherwise.}\end{cases}
		\end{align*}
		Hence $\Ind(P)$ is equivalent to the poset of \emph{ideals} of $P$, i.e.\ non-empty, directed, downward closed subsets $J \subseteq P$, ordered by inclusion, via $S \mapsto {\downarrow}S := \{p \in P \mid p \leq s \text{ for some } s \in S\}$.
	\end{lemma}
	\prooffor{lem:Ind_Of_Poset}
	\begin{proof}
		By \Cref{thm:Filtered_Posets_Cofinal}, every object of $\Ind(P)$ is a colimit of representables indexed by a filtered poset $I$, i.e.\ of the form $\quotcolim_{i \in I} p_i$ for a functor $p\colon I \to P$. The image $S = p(I)$ is a directed subset, and the map $I \to S$ is cofinal by \Cref{lem:Cofinality_Posets} (for $s \in S$, the set $\{i \mid s \leq p_i\}$ is non-empty and upward closed in the filtered poset $I$, hence filtered). This proves the first claim. The formula for mapping spaces follows from \Cref{constr:Ind_Large}: each $\Map_P(s,t)$ is either empty or contractible, filtered colimits of such spaces are again empty or contractible (and contractible if and only if one of the terms is), and a limit of such spaces is contractible if and only if all terms are. In particular, all mapping spaces in $\Ind(P)$ are empty or contractible, so $\Ind(P)$ is a poset. Finally, the condition in the formula only depends on ${\downarrow}S$ and ${\downarrow}T$: it says precisely that ${\downarrow}S \subseteq {\downarrow}T$.
	\end{proof}

	Suppose now that $P$ admits directed suprema. Then the colimit functor $\colim\colon \Ind(P) \to P$ sends a directed subset $S$ to $\sup S$, and the Yoneda embedding sends $b$ to ${\downarrow}b$.

	\begin{lemma}
		\label[lemma]{lem:Posets_Accessible}
		Let $P$ be a small partially ordered set admitting directed suprema. Then $P$ is accessible.
	\end{lemma}
	\prooffor{lem:Posets_Accessible}
	\begin{proof}
		Let $\kappa$ be a regular cardinal with $\kappa > |P|$. Let $D \subseteq P$ be a $\kappa$-directed subset. Since $|D| < \kappa$, the set $D$ itself has an upper bound in $D$ (\Cref{lem:Kappa_Filtered_Posets}), i.e.\ $D$ has a largest element. It follows that every $\kappa$-filtered diagram in $P$ has a cofinal constant subdiagram, so that $\kappa$-filtered colimits in $P$ are ``trivial'': every object of $P$ is $\kappa$-compact, and the Yoneda embedding $P \to \Ind_{\kappa}(P)$ is an equivalence, as every $\kappa$-filtered colimit of representable presheaves on $P$ is representable. Hence $P \simeq \Ind_{\kappa}(P^{\kappa})$ is $\kappa$-accessible.
	\end{proof}

	\begin{definition}
		\label[definition]{def:Continuous_Poset}
		A partially ordered set $P$ is called \emph{continuous} if it admits directed suprema and the functor $\colim = \sup\colon \Ind(P) \to P$ admits a left adjoint $\hat{y}\colon P \to \Ind(P)$.
	\end{definition}

	By \Cref{lem:Posets_Accessible}, a small poset is continuous if and only if it is compactly assembled. Unwinding the definitions using \Cref{lem:Ind_Of_Poset}, the existence of the left adjoint $\hat{y}$ means: for every $b \in P$ there is a directed subset $\hat{y}(b) \subseteq P$ such that for every directed subset $T \subseteq P$,
	\begin{equation}
		\label{eq:Yhat_Poset}
		b \leq \sup T \iff \text{for all } s \in \hat{y}(b) \text{ there exists } t \in T \text{ with } s \leq t.
	\end{equation}

	\begin{definition}
		\label[definition]{def:Way_Below}
		Let $P$ be a poset admitting directed suprema, and let $a,b \in P$. We say that $a$ is \emph{way below} $b$, written $a \ll b$, if for every directed subset $T \subseteq P$ with $b \leq \sup T$ there exists $t \in T$ with $a \leq t$. We write
		\[
			P_{\ll b} := \{a \in P \mid a \ll b\}.
		\]
	\end{definition}

	\begin{remark}
		\label[remark]{rem:Way_Below_Basics}
		\leavevmode
		\begin{enumerate}
			\item If $a \ll b$, then $a \leq b$: apply the definition to $T = \{b\}$.
			\item If $a' \leq a \ll b \leq b'$, then $a' \ll b'$.
			\item An element $b$ is compact in $P$ (i.e.\ $\Map_P(b,-)$ preserves filtered colimits) if and only if $b \ll b$.
			\item If $P$ is continuous, then $\hat{y}(b) \subseteq P_{\ll b}$ for every $b$. Indeed, let $s \in \hat{y}(b)$ and let $T$ be directed with $b \leq \sup T$. By \eqref{eq:Yhat_Poset}, there exists $t \in T$ with $s \leq t$, so $s \ll b$.
		\end{enumerate}
	\end{remark}

	\begin{proposition}
		\label[proposition]{prop:Continuous_Poset_Characterization}
		Let $P$ be a poset admitting directed suprema. Then $P$ is continuous if and only if for every $b \in P$ the subset $P_{\ll b}$ is directed and $\sup P_{\ll b} = b$. In this case, $\hat{y}(b) = P_{\ll b}$, viewed as an object of $\Ind(P)$.
	\end{proposition}
	\prooffor{prop:Continuous_Poset_Characterization}
	\begin{proof}
		Suppose that each $P_{\ll b}$ is directed with supremum $b$. We claim that $b \mapsto P_{\ll b}$ satisfies \eqref{eq:Yhat_Poset}, and hence defines a left adjoint to $\sup$ (a monotone map, by \Cref{rem:Way_Below_Basics}(2)). If $b \leq \sup T$, then by definition every $a \ll b$ satisfies $a \leq t$ for some $t \in T$. Conversely, if every $a \in P_{\ll b}$ is bounded by some element of $T$, then $b = \sup P_{\ll b} \leq \sup T$.

		Conversely, suppose that $P$ is continuous. We claim that
		\[
			P_{\ll b} = \{a \in P \mid a \leq s \text{ for some } s \in \hat{y}(b)\} = {\downarrow}\hat{y}(b).
		\]
		For ``$\supseteq$'', note that $\hat{y}(b) \subseteq P_{\ll b}$ by \Cref{rem:Way_Below_Basics}(4), and $P_{\ll b}$ is downward closed by \Cref{rem:Way_Below_Basics}(2). For ``$\subseteq$'', note that $\hat{y}$ is fully faithful (\Cref{lem:Adjoint_Triples}), so the counit $\sup\hat{y}(b) \leq b$ is an equality, i.e.\ $b = \sup \hat{y}(b)$. Applying the definition of $a \ll b$ to the directed set $T = \hat{y}(b)$ gives $s \in \hat{y}(b)$ with $a \leq s$. Since $\hat{y}(b)$ is directed, so is ${\downarrow}\hat{y}(b) = P_{\ll b}$, and $\sup P_{\ll b} = \sup \hat{y}(b) = b$. Finally, $P_{\ll b}$ and $\hat{y}(b)$ define the same object of $\Ind(P)$ by \Cref{lem:Ind_Of_Poset}.
	\end{proof}

	\begin{example}
		\label[example]{ex:Continuous_Poset_Reals}
		The poset $P = \R \cup \{\infty\}$ (with its usual order) admits directed suprema: the supremum of a non-empty subset is its usual supremum if the subset is bounded, and $\infty$ otherwise. We claim that $a \ll b$ if and only if $a < b$. If $a < b$ and $T$ is directed with $b \leq \sup T$, then $\sup T > a$, so some $t \in T$ satisfies $t > a$. Conversely, suppose $a \ll b$; then $a \leq b$ by \Cref{rem:Way_Below_Basics}(1). If $a = b$, consider the directed set $T = \{b - 1/n \mid n \geq 1\}$ (for $b \in \R$) resp.\ $T = \N$ (for $b = \infty$): it has supremum $b$, but no element $t \in T$ satisfies $b \leq t$. Hence $a < b$. It follows that $P_{\ll b} = \{a \mid a < b\}$ is directed with supremum $b$, so $P$ is continuous by \Cref{prop:Continuous_Poset_Characterization}, with $\hat{y}(b) = \quotcolim_{a < b} a$. Note that $P$ has no compact elements.
	\end{example}

	\begin{example}
		\label[example]{ex:Continuous_Poset_Opens}
		Let $X$ be a locally compact Hausdorff space and let $\Open(X)$ be the poset of open subsets of $X$, ordered by inclusion. Directed suprema in $\Open(X)$ are unions. We claim that
		\[
			U \ll V \iff \overline{U} \text{ is compact and } \overline{U} \subseteq V,
		\]
		in which case we write $U \Subset V$. Suppose first that $\overline{U}$ is compact and $\overline{U} \subseteq V$, and let $(W_i)_{i \in I}$ be a directed family of open subsets with $V \subseteq \bigcup_i W_i$. Then the $W_i$ cover the compact set $\overline{U}$, so finitely many of them do, and since the family is directed, a single $W_i$ contains $\overline{U} \supseteq U$. Conversely, let $U \ll V$. Since $X$ is locally compact Hausdorff, every point $x \in V$ has an open neighbourhood $W$ with $\overline{W}$ compact and $\overline{W} \subseteq V$. Hence $V$ is the union of the family $\{W \in \Open(X) \mid W \Subset V\}$, which is directed since $W_1 \cup W_2 \Subset V$ if $W_1,W_2 \Subset V$. By assumption, $U \subseteq W$ for some $W \Subset V$, and then $\overline{U} \subseteq \overline{W}$ is a closed subset of a compact set, hence compact, and $\overline{U} \subseteq \overline{W} \subseteq V$.

		The argument shows moreover that $\Open(X)_{\ll V}$ is directed with union $V$, so $\Open(X)$ is a continuous poset by \Cref{prop:Continuous_Poset_Characterization}, with $\hat{y}(V) = \quotcolim_{U \Subset V} U$. Its compact elements are the open subsets $U$ with $U \Subset U$, i.e.\ the compact open subsets of $X$.
	\end{example}

	\begin{remark}
		\label[remark]{rem:Continuous_Posets_Sheaves}
		If $X$ is a locally compact Hausdorff space, then the morphism $h_U \to h_V$ of representable sheaves on $X$ is a compact morphism in $\Shv(X;\Sp)$ whenever $U \Subset V$ (Mathew, Lecture 3). Since every open $V$ is the union of a sequence $U_0 \Subset U_1 \Subset \cdots$ if $X$ is second countable, this suggests (and, suitably generalized, proves) that $\Shv(X;\Sp)$ is generated by sequential colimits along compact morphisms and hence dualizable, cf.\ \Cref{rem:Omega_1_Compact_Compact_Morphisms}. We will return to this in \Cref{chap:Sheaves}.
	\end{remark}

	\section{Rigid symmetric monoidal categories}
	\label[section]{sec:Rigid_Categories}

	% Source: Talk 3 (Max Glaser, 27.10.2026); Clausen, Lecture 3 (19.07.2023); literature (Gaitsgory--Rozenblyum, Ramzi, Efimov), proofs partly completed with generative AI (Claude)
	A rich source of dualizable categories are \emph{rigid} symmetric monoidal categories, introduced by Gaitsgory--Rozenblyum \cite[Chapter~1, Section~9]{GaitsgoryRozenbluym2017StudyDAG}. We have already met them in \Cref{def:Rigid_Category}; we now give a more detailed treatment.

	\begin{definition}[{\cite[Chapter~1, Definition~9.1.2]{GaitsgoryRozenbluym2017StudyDAG}}]
		\label[definition]{def:Rigid_Category_Detailed}
		Let $C \in \CAlg(\PrL_{\st})$ be a presentably symmetric monoidal stable category, with unit $\unit_C$ and multiplication $\mu\colon C \otimes C \to C$. Then $C$ is called \emph{rigid} if
		\begin{enumerate}
			\item the unit $\unit_C$ is compact in $C$, or equivalently, the unit functor $u\colon \Sp \to C$, $\S \mapsto \unit_C$, is strongly continuous;
			\item the right adjoint $\mu^R\colon C \to C \otimes C$ preserves colimits, and its canonical lax $C \otimes C$-linear structure is strict; equivalently, $\mu$ admits a right adjoint in $\Mod_{C \otimes C}(\PrL_{\st})$.
		\end{enumerate}
	\end{definition}

	\begin{remark}
		\label[remark]{rem:Equivalent_Rigidity_Definitions}
		The equivalence of the two formulations in (2) is a general fact about adjunctions internal to module categories: a right adjoint of a $C \otimes C$-linear functor inherits a lax $C \otimes C$-linear structure, and it is a right adjoint in $\Mod_{C \otimes C}(\PrL_{\st})$ if and only if it preserves colimits and this lax structure is strict. \Cref{def:Rigid_Category_Detailed} is the definition of Gaitsgory--Rozenblyum, and it agrees with Ramzi's definition \cite[Definitions~4.5 and~4.36]{Ramzi_Locally_Rigid} (Ramzi additionally requires $C$ to be dualizable, which is automatic by \Cref{prop:Rigid_Self_Dual_Detailed}). It is also equivalent to \Cref{def:Rigid_Category}, which asks that for every finite set $I$ the multiplication $\mu_I\colon C^{\otimes I} \to C$ admits a right adjoint in $\Mod_{C^{\otimes I}}(\PrL_{\st})$. Indeed, for $I = \emptyset$ the map $\mu_{\emptyset}$ is the unit functor $u\colon \Sp \to C$, and $\Mod_{\Sp}(\PrL_{\st}) = \PrL_{\st}$, so this case is condition (1); the case $|I| = 1$ is vacuous, and the case $|I| = 2$ is condition (2). Conversely, if (2) holds, then by induction $\mu_{I \sqcup \{*\}} = \mu \circ (\mu_I \otimes \id_C)$ is a composite of two functors admitting right adjoints in $\Mod_{C^{\otimes I} \otimes C}(\PrL_{\st})$: for $\mu_I \otimes \id_C$ this follows by tensoring the adjunction for $\mu_I$ with $C$, and for $\mu$ by restricting the $C \otimes C$-linear adjunction along the algebra map $\mu_I \otimes \id_C\colon C^{\otimes I} \otimes C \to C \otimes C$.
	\end{remark}

	\begin{proposition}[{\cite[Chapter~1, Section~9.2.1]{GaitsgoryRozenbluym2017StudyDAG}}]
		\label[proposition]{prop:Rigid_Self_Dual_Detailed}
		Let $C$ be rigid. Then $C$ is self-dual in $\PrL_{\st}$, with evaluation and coevaluation given by
		\[
			\ev\colon C \otimes C \xrightarrow{\mu} C \xrightarrow{\Gamma := \hom_C(\unit_C,-)} \Sp, \qquad \coev\colon \Sp \xrightarrow{u} C \xrightarrow{\mu^R} C \otimes C.
		\]
		In particular, $C$ is dualizable.
	\end{proposition}
	\prooffor{prop:Rigid_Self_Dual_Detailed}
	\begin{proof}
		Note first that both functors preserve colimits: $\Gamma$ preserves colimits since $\unit_C$ is compact (and $C$ is stable), and $\mu^R$ preserves colimits by assumption. Moreover, $\Gamma$ is right adjoint to $u$. Write $\eta := \mu^R(\unit_C) \in C \otimes C$, so that $\coev$ is the colimit-preserving functor $\Sp \to C \otimes C$ sending $\S$ to $\eta$. We verify the first triangle identity, i.e.\ that the composite
		\[
			C \simeq \Sp \otimes C \xrightarrow{\coev \otimes C} C \otimes C \otimes C \xrightarrow{C \otimes \ev} C \otimes \Sp \simeq C
		\]
		is equivalent to the identity; the second one follows by symmetry. Since $C \otimes \ev = (C \otimes \Gamma) \circ (C \otimes \mu)$, the composite sends $c \in C$ to
		\[
			(C \otimes \Gamma)\big((C \otimes \mu)(\eta \boxtimes c)\big) \simeq (C \otimes \Gamma)\big(\eta \cdot (\unit_C \boxtimes c)\big),
		\]
		where $\cdot$ denotes the multiplication in $C \otimes C$. Now we use the linearity of $\mu^R$: for $a \in C \otimes C$ and $x \in C$ we have $\mu^R(\mu(a) \otimes x) \simeq a \cdot \mu^R(x)$. Taking $a = \unit_C \boxtimes c$ and $x = \unit_C$ gives
		\[
			\eta \cdot (\unit_C \boxtimes c) = (\unit_C \boxtimes c) \cdot \mu^R(\unit_C) \simeq \mu^R(c).
		\]
		Hence the composite is $(C \otimes \Gamma) \circ \mu^R\colon C \to C$, naturally in $c$ (as all functors involved preserve colimits, it suffices to identify the functors, which the above does naturally in $c$). Since $C \otimes - $ is a functor of $(\infty,2)$-categories, $C \otimes \Gamma$ is right adjoint to $C \otimes u\colon C \to C \otimes C$, $x \mapsto x \boxtimes \unit_C$. Hence $(C \otimes \Gamma) \circ \mu^R$ is right adjoint to $\mu \circ (C \otimes u) \simeq \id_C$, and is therefore equivalent to the identity.
	\end{proof}

	Next, we introduce trace-class morphisms, following Clausen--Scholze \cite{Scholze_Analytic_Geometry}.

	\begin{definition}
		\label[definition]{def:Trace_Class}
		Let $C$ be a symmetric monoidal category. A morphism $f\colon x \to y$ in $C$ is called \emph{trace-class} if there exist an object $z \in C$ and maps $\alpha\colon \unit_C \to y \otimes z$ and $\beta\colon z \otimes x \to \unit_C$ such that $f$ is homotopic to the composite
		\[
			x \simeq \unit_C \otimes x \xrightarrow{\alpha \otimes x} y \otimes z \otimes x \xrightarrow{y \otimes \beta} y \otimes \unit_C \simeq y.
		\]
	\end{definition}

	\begin{remark}
		\label[remark]{rem:Trace_Class_Basics}
		\leavevmode
		\begin{enumerate}
			\item If $C$ is closed, write $x^{\vee} := \iHom(x,\unit_C)$ for the \emph{predual} of $x$. Then $f$ is trace-class if and only if the map $\unit_C \to \iHom(x,y)$ classifying $f$ factors through the canonical map $y \otimes x^{\vee} \to \iHom(x,y)$: given $(z,\alpha,\beta)$ as above, $\beta$ is adjoint to a map $z \to x^{\vee}$, and composing $\alpha$ with it gives the desired factorization; conversely, given a factorization $\unit_C \to y \otimes x^{\vee}$, take $z = x^{\vee}$ and $\beta$ the evaluation map.
			\item Trace-class morphisms form a two-sided ideal: if $f$ is trace-class with witness $(z,\alpha,\beta)$, then $hfk$ is trace-class with witness $(z,(h \otimes z)\alpha,\beta(z \otimes k))$.
			\item If $C$ is stable and idempotent complete (for instance presentable and stable), then the identity of $x$ is trace-class if and only if $x$ is dualizable, see \Cref{lem:Identity_Trace_Class} below.
			\item A symmetric monoidal functor preserves trace-class morphisms, since it preserves the data $(z,\alpha,\beta)$.
		\end{enumerate}
	\end{remark}

	\begin{lemma}
		\label[lemma]{lem:Identity_Trace_Class}
		Let $C$ be a symmetric monoidal category whose homotopy category is idempotent complete; for instance, $C$ stable and idempotent complete \cite[Lemma~1.2.4.6]{lurie2016HA}. Then an object $x \in C$ is dualizable if and only if $\id_x$ is trace-class.
	\end{lemma}
	\prooffor{lem:Identity_Trace_Class}
	\begin{proof}
		If $x$ is dualizable with dual $x^{\vee}$, then the coevaluation and evaluation are a witness for $\id_x$. Conversely, let $(z,\alpha,\beta)$ be a witness for $\id_x$, so that $\alpha\colon \unit \to x \otimes z$ and $\beta\colon z \otimes x \to \unit$ satisfy the first triangle identity $(x \otimes \beta)(\alpha \otimes x) \simeq \id_x$, but not necessarily the second one. Consider
		\[
			e := (\beta \otimes z)(z \otimes \alpha)\colon z \to z \otimes x \otimes z \to z.
		\]
		Inserting the first triangle identity in the middle of $e \circ e$ shows that $e \circ e \simeq e$, and similarly $\beta(e \otimes x) \simeq \beta$ and $(x \otimes e)\alpha \simeq \alpha$. As the homotopy category is idempotent complete, $e$ splits: $e \simeq sr$ with $r\colon z \to z'$, $s\colon z' \to z$ and $rs \simeq \id_{z'}$. Put $\alpha' := (x \otimes r)\alpha$ and $\beta' := \beta(s \otimes x)$. Then
		\[
			(x \otimes \beta')(\alpha' \otimes x) \simeq (x \otimes \beta)(x \otimes e \otimes x)(\alpha \otimes x) \simeq (x \otimes \beta)(\alpha \otimes x) \simeq \id_x
		\]
		and $(\beta' \otimes z')(z' \otimes \alpha') \simeq r e s \simeq rsrs \simeq \id_{z'}$. So $(\alpha',\beta')$ exhibits $z'$ as a dual of $x$; since dualizability can be checked in the homotopy category, $x$ is dualizable.
	\end{proof}

	\begin{lemma}
		\label[lemma]{lem:Trace_Class_Compact}
		Let $C \in \CAlg(\PrL_{\st})$ be such that $\unit_C$ is compact. Then every trace-class morphism in $C$ is compact in the sense of \Cref{def:Compact_Morphism}.
	\end{lemma}
	\prooffor{lem:Trace_Class_Compact}
	\begin{proof}
		Let $f\colon x \to y$ be trace-class with witness $(z,\alpha,\beta)$. We construct a factorization of $f^*\colon \Map_C(y,\colim(-)) \to \Map_C(x,\colim(-))$ through $\Map_{\Ind(C)}(y(x),-)$, naturally in $W = \quotcolim_j w_j \in \Ind(C)$. Given $g\colon y \to \colim_j w_j$, consider
		\[
			(g \otimes z) \circ \alpha\colon \unit_C \to (\colim_j w_j) \otimes z \simeq \colim_j (w_j \otimes z).
		\]
		Since $\unit_C$ is compact, this defines a point of $\colim_j \Map_C(\unit_C,w_j \otimes z)$, and postcomposing with the natural map
		\[
			\colim_j \Map_C(\unit_C,w_j \otimes z) \to \colim_j \Map_C(x,w_j), \qquad \gamma \mapsto (w_j \otimes \beta) \circ (\gamma \otimes x),
		\]
		we obtain a point of $\colim_j \Map_C(x,w_j) = \Map_{\Ind(C)}(y(x),W)$. All these maps are natural in $W$, and the image of this point in $\Map_C(x,\colim_j w_j)$ is $(\colim_j w_j \otimes \beta) \circ ((g \otimes z)\alpha \otimes x) \simeq g \circ (y \otimes \beta) \circ (\alpha \otimes x) \simeq gf$.
	\end{proof}

	The converse holds for rigid categories; this is the main structural result about rigid categories. It was formulated in Clausen's lectures as the \emph{definition} of rigidity for dualizable symmetric monoidal categories.

	\begin{theorem}[{Ramzi \cite[Corollaries~2.6, 4.55 and~4.60, Example~2.9]{Ramzi_Locally_Rigid}; see also \cite[Propositions~1.6 and~1.7]{Efimov_Inverse_Limits}}]
		\label[theorem]{thm:Rigid_Trace_Class}
		Let $C \in \CAlg(\PrL_{\st})$. Then the following are equivalent:
		\begin{enumerate}
			\item $C$ is rigid;
			\item $\unit_C$ is compact, and $C$ is generated under colimits by objects of the form $\colim(x_0 \to x_1 \to x_2 \to \cdots)$ in which each map $x_n \to x_{n+1}$ is trace-class;
			\item $C$ is dualizable, and a morphism in $C$ is compact if and only if it is trace-class.
		\end{enumerate}
	\end{theorem}

	We will not prove this theorem. Note however that some implications are formal given what we know: (1) $\implies$ the first half of (3) is \Cref{prop:Rigid_Self_Dual_Detailed}, and (3) $\implies$ (2) follows from \Cref{lem:Compact_Morphisms_Yhat}(2), since $\id_{\unit_C}$ is trace-class and hence compact, and \Cref{rem:Omega_1_Compact_Compact_Morphisms}, which shows that a dualizable category is generated by sequential colimits along compact morphisms. The hard part is that in a rigid category every compact morphism is trace-class \cite[Corollary~4.55, Example~2.9 and Corollary~2.6]{Ramzi_Locally_Rigid}, and that (2) implies (1) \cite[Corollary~4.60]{Ramzi_Locally_Rigid}. Note that the dualizability hypothesis in (3) was not stated in the talk; we do not know whether it can be dropped.

	\begin{corollary}
		\label[corollary]{cor:Rigid_Functors_Strongly_Continuous}
		Let $F\colon C \to D$ be a colimit-preserving symmetric monoidal functor between rigid categories. Then $F$ is strongly continuous. Consequently, the category $\mathrm{Rig}$ of rigid symmetric monoidal categories and colimit-preserving symmetric monoidal functors is a full subcategory of $\CAlg(\Pr^{\dual})$.
	\end{corollary}
	\prooffor{cor:Rigid_Functors_Strongly_Continuous}
	\begin{proof}
		By \Cref{thm:Rigid_Trace_Class}, the compact morphisms in $C$ are trace-class. Since $F$ is symmetric monoidal, it preserves trace-class morphisms (\Cref{rem:Trace_Class_Basics}(4)), and trace-class morphisms in $D$ are compact by \Cref{lem:Trace_Class_Compact}. So $F$ preserves compact morphisms, and hence is strongly continuous by \Cref{prop:Strongly_Continuous_Compact_Morphisms}.
	\end{proof}

	\begin{remark}
		\label[remark]{rem:Strong_Continuity_From_Rigid_Source}
		The argument shows more generally that a colimit-preserving symmetric monoidal functor $F\colon C \to D$ with $C$ rigid and $D$ dualizable with compact unit is strongly continuous. In fact, the dualizability hypothesis on $D$ can be removed \cite[Proposition~1.9]{Efimov_Inverse_Limits}.
	\end{remark}

	\begin{example}
		\label[example]{ex:Rigid_Compactly_Generated}
		Let $C \in \CAlg(\PrL_{\st})$ be compactly generated. Then $C$ is rigid if and only if $\unit_C$ is compact and every compact object of $C$ is dualizable, if and only if the compact objects of $C$ are precisely the dualizable objects \cite[Chapter~1, Lemma~9.1.5 and Corollary~9.1.7]{GaitsgoryRozenbluym2017StudyDAG}. Indeed, if $\unit_C$ is compact, then dualizable objects are compact (as $\Map(x,-) \simeq \Map(\unit_C,x^{\vee} \otimes -)$), and conversely if all compact objects are dualizable, then $C = \Ind(C^{\omega})$ is generated by objects whose identity is trace-class, so (2) of \Cref{thm:Rigid_Trace_Class} holds. Examples include $\D(R)$ for a commutative ring $R$, more generally $\D_{\mathrm{qc}}(X)$ for a qcqs scheme $X$, and $\Sp$ itself.
	\end{example}

	\begin{example}
		\label[example]{ex:Shv_Compact_Hausdorff_Rigid}
		Let $X$ be a compact Hausdorff space. Then $\Shv(X;\Sp)$ is rigid; this can be seen as a manifestation of Verdier duality, see \cite[Example~1.4]{Efimov_Inverse_Limits}. For instance, the unit (the constant sheaf) is compact since global sections $\Gamma(X;-)$ commute with filtered colimits for compact $X$ \cite[Corollary~7.3.4.12]{lurie2009HTT}. In particular, $\Shv(X;\Sp)$ is dualizable and self-dual. For every presentable stable category $D$ we have $\Shv(X;D) \simeq \Shv(X;\Sp) \otimes D$ \cite[Proposition~2.8.7]{Krause_Nikolaus_Putzstuck_Sheaves}, so $\Shv(X;D)$ is dualizable whenever $D$ is. We will study these categories in detail in \Cref{chap:Sheaves}.
	\end{example}

	\begin{example}
		\label[example]{ex:Locally_Rigid_Sheaf_Categories}
		If $X$ is merely locally compact Hausdorff and not compact, then $\Shv(X;\Sp)$ is not rigid, since the unit is not compact. It is still dualizable, and even \emph{locally rigid} in the sense of Ramzi \cite[Definition~4.5 and Example~4.10]{Ramzi_Locally_Rigid} (see also \cite[Definition~1.36 and Example~1.38]{Efimov_Inverse_Limits} and \cite[Example~4.2.9, Corollary~4.4.13]{Krause_Nikolaus_Putzstuck_Sheaves}): it satisfies condition (2) of \Cref{def:Rigid_Category_Detailed}. Indeed, for $j\colon X \hookrightarrow X^+$ the one-point compactification, $j_!$ identifies $\Shv(X;\Sp)$ with a smashing ideal in the rigid category $\Shv(X^+;\Sp)$.
	\end{example}

	\section{Nuclear modules and condensed mathematics}
	\label[section]{sec:Nuclear_Modules}

	% Source: Talk 3 (Max Glaser, 27.10.2026); literature (Clausen--Scholze, Efimov, Ramzi, Andreychev), drafted with generative AI (Claude)
	The most important examples of dualizable categories which are not compactly generated arise in condensed mathematics, as categories of \emph{nuclear modules}. These are extracted from large compactly generated categories of ``complete'' modules, which are not rigid, via the following general construction due to Clausen and Scholze \cite{Scholze_Analytic_Geometry}.

	Throughout this section, let $C \in \CAlg(\PrL_{\st})$ be compactly generated with compact unit. We write $x^{\vee} = \iHom(x,\unit_C)$ for the predual of $x$.

	\begin{definition}[{Clausen--Scholze, \cite[Definition~8.5]{Clausen_Scholze_Complex_Geometry}; cf.\ \cite[Definitions~13.10 and~13.12]{Scholze_Analytic_Geometry}}]
		\label[definition]{def:Nuclear}
		An object $M \in C$ is called \emph{nuclear} if every morphism $P \to M$ from a compact object $P \in C^{\omega}$ is trace-class. We denote by $\Nuc(C) \subseteq C$ the full subcategory of nuclear objects. An object $M$ is called \emph{basic nuclear} if $M \simeq \colim(P_0 \to P_1 \to P_2 \to \cdots)$ with each $P_n$ compact and each map $P_n \to P_{n+1}$ trace-class.
	\end{definition}

	Note that a compact object $P$ is nuclear if and only if $\id_P$ is trace-class, i.e.\ if and only if $P$ is dualizable (\Cref{rem:Trace_Class_Basics}(3)). So the nuclear objects should be thought of as the objects which are ``locally dualizable''.

	\begin{proposition}[Clausen--Scholze; see {\cite[Theorem~8.6]{Clausen_Scholze_Complex_Geometry} and \cite[Propositions~1.27 and~1.28]{Efimov_Inverse_Limits}}]
		\label[proposition]{prop:Nuclear_Basics}
		\leavevmode
		\begin{enumerate}
			\item For $M \in C$, the following are equivalent:
			\begin{enumerate}
				\item $M$ is nuclear;
				\item for every compact $P$, the map $\Map_C(\unit_C,P^{\vee} \otimes M) \to \Map_C(P,M)$ is an equivalence;
				\item for every compact $P$, the map $P^{\vee} \otimes M \to \iHom(P,M)$ is an equivalence;
				\item $M \simeq \colim_{i \in I} M_i$ for a directed poset $I$ and compact objects $M_i$, such that for every $i$ there is $j \geq i$ with $M_i \to M_j$ trace-class.
			\end{enumerate}
			\item The subcategory $\Nuc(C) \subseteq C$ is closed under colimits. It is $\omega_1$-presentable, and its $\omega_1$-compact objects are precisely the basic nuclear objects.
		\end{enumerate}
	\end{proposition}

	In general, $\Nuc(C)$ need not be rigid, or even dualizable. The following criterion is the one that is used in practice.

	\begin{proposition}[{\cite[Propositions~1.23 and~1.29]{Efimov_Inverse_Limits}, \cite[Construction~4.79 and Theorem~4.81]{Ramzi_Locally_Rigid}}]
		\label[proposition]{prop:Nuclear_Rigid}
		Suppose that every trace-class morphism between compact objects of $C$ is a composite of two trace-class morphisms. Then $\Nuc(C)$ is a rigid symmetric monoidal category. More precisely, $\Nuc(C)$ is equivalent to the \emph{rigidification} $C^{\rig}$ of $C$, i.e.\ the value at $C$ of the right adjoint to the inclusion of rigid categories into $\CAlg(\PrL_{\st})$.
	\end{proposition}
	\prooffor{prop:Nuclear_Rigid}
	\begin{proof}[Sketch]
		By \cite[Proposition~1.23]{Efimov_Inverse_Limits}, based on \cite[Theorem~4.81 and Corollary~4.75]{Ramzi_Locally_Rigid}, the rigidification $C^{\rig} \subseteq C$ is the full subcategory generated under colimits by colimits $\colim_{a \in \Q} F(a)$ of diagrams $F\colon (\Q,\leq) \to C^{\omega}$ in which every map $F(a) \to F(b)$, $a < b$, is trace-class. Such colimits are nuclear by \Cref{prop:Nuclear_Basics}(1)(d), so $C^{\rig} \subseteq \Nuc(C)$. Conversely, the assumption allows one to interpolate any sequence $P_0 \to P_1 \to \cdots$ of trace-class maps between compact objects to a $\Q$-indexed diagram of trace-class maps (by repeatedly factoring the maps), so the basic nuclear objects lie in $C^{\rig}$. By \Cref{prop:Nuclear_Basics}(2), these generate $\Nuc(C)$.
	\end{proof}

	\begin{remark}
		\label[remark]{rem:Rigidification_Via_Idempotent_Ideals}
		Compare this with the construction of $(C,S)^{\dual}$ in Clausen's lectures: given a class $S$ of morphisms forming an ideal with $S = S^2$ (and some further closure properties), one obtains a dualizable category generated by colimits along morphisms in $S$. For $S$ the trace-class morphisms of $C$, this gives $C^{\rig}$. The condition $S = S^2$ in \Cref{prop:Nuclear_Rigid} is what allows one to work with sequential colimits instead of $\Q$-indexed ones, cf.\ also the construction of $C_S$ in Lecture~7 of the Bonn workshop.
	\end{remark}

	\subsection*{Examples from condensed mathematics}

	We now describe the main examples, without proofs. We refer to \cite{Scholze_Analytic_Geometry} for the solid and liquid theory, to \cite{Clausen_Scholze_Complex_Geometry} for the gaseous theory, and to \cite{Andreychev_Thesis} and \cite[Section~7.2]{Efimov_Inverse_Limits} for the comparison with Efimov's theory.

	\begin{example}[Solid modules and nuclear modules over formal schemes]
		\label[example]{ex:Nuclear_Solid}
		Let $\D(\mathrm{Solid}_{\Z})$ be the derived category of solid abelian groups. This is a compactly generated symmetric monoidal category, with compact projective generator $\prod_{\N}\Z$ and compact unit $\Z$. It is far from rigid: the compact object $\prod_{\N}\Z$ is not dualizable (its predual is $\bigoplus_{\N}\Z$). In fact, there is a symmetric monoidal equivalence
		\[
			\D(\mathrm{Solid}_{\Z}) \simeq \Ind\big(\D^b(\mathrm{Proj}^{\omega_1}\text{-}\Z)\catop\big),
		\]
		where $\D^b(\mathrm{Proj}^{\omega_1}\text{-}\Z) \subseteq \D(\Z)$ is the stable subcategory generated by $\bigoplus_{\N}\Z$, under which $\prod_{\N}\Z$ corresponds to $(\bigoplus_{\N}\Z)\catop$ \cite[Section~7.2]{Efimov_Inverse_Limits}.

		More generally, let $R$ be a commutative ring and $I = (a_1,\dots,a_m) \subseteq R$ a finitely generated ideal, and consider the $I$-adic completion $R^{\wedge}_I$ as an analytic ring with the induced solid structure. Its category of modules $\D(\mathrm{Solid}_{(R^{\wedge}_I,\Z)}) = \Mod_{R^{\wedge}_I}(\D(\mathrm{Solid}_{\Z}))$ is compactly generated by $R^{\wedge}_I \otimes^{\square} \prod_{\N}\Z$, with compact unit. The category of \emph{nuclear solid modules} is
		\[
			\Nuc^{CS}(R^{\wedge}_I) := \Nuc\big(\D(\mathrm{Solid}_{(R^{\wedge}_I,\Z)})\big),
		\]
		which does not depend on the choice of analytic ring structure on $R^{\wedge}_I$. (For $R = \Z$ and $I = 0$ one recovers $\D(\Z)$: an object of $\D(\mathrm{Solid}_{\Z})$ is nuclear if and only if it is discrete \cite[Lecture~VIII, Exercise~1]{Clausen_Scholze_Complex_Geometry}.) One has $\Nuc^{CS}(R^{\wedge}_I) \simeq \D(\mathrm{Solid}_{(R^{\wedge}_I,\Z)})^{\rig}$. The key technical input is that the preduals of compact objects are nuclear, which goes back to Andreychev \cite[Lemma~3.8 and Korollar~3.14]{Andreychev_Thesis}; see \cite[Propositions~7.4 and~7.5]{Efimov_Inverse_Limits} (Efimov deduces the rigidity from this via \cite[Propositions~1.33 and~1.34]{Efimov_Inverse_Limits} rather than via \Cref{prop:Nuclear_Rigid}). This boils down to the formula $\Lambda_I(\bigoplus_{\N}R) \otimes \Lambda_I(\bigoplus_{\N}R) \simeq \Lambda_I(\bigoplus_{\N \times \N}R)$ for a suitable completion functor $\Lambda_I$). Hence:
		\begin{itemize}
			\item $\Nuc^{CS}(R^{\wedge}_I)$ is rigid, hence dualizable, but in general not compactly generated;
			\item its compact objects are the perfect complexes over $R^{\wedge}_I$;
			\item its continuous K-theory is $K^{\cont}(\Nuc^{CS}(R^{\wedge}_I)) \simeq \lim_n K(R \quot (a_1^n,\dots,a_m^n))$, which is $\lim_n K(R/I^n)$ if $R$ is Noetherian (Efimov \cite[Corollary~7.10]{Efimov_Inverse_Limits}, as conjectured by Clausen--Scholze), cf.\ \Cref{thm:Kcont_Nuclear}.
		\end{itemize}
		Efimov also introduced a variant $\Nuc(R^{\wedge}_I)$ \cite[Definition~3.28]{Efimov_Inverse_Limits}, which can be described as $\lim^{\dual}_n \D(R \quot (a_1^n,\dots,a_m^n))$, the limit computed in $\Pr^{\dual}$ \cite[Proposition~5.23]{Efimov_Inverse_Limits}. It receives a fully faithful strongly continuous functor from $\Nuc^{CS}(R^{\wedge}_I)$ \cite[Corollary~7.6]{Efimov_Inverse_Limits}, which induces an equivalence on compact objects, both sides having compact objects $\Perf(R^{\wedge}_I)$ \cite[Proposition~3.29]{Efimov_Inverse_Limits}. It can equivalently be described as the rigidification $\D_{I\text{-}\mathrm{compl}}(R)^{\rig}$ \cite[Proposition~4.3]{Efimov_Inverse_Limits} of the category of derived $I$-complete $R$-modules; the latter is itself not rigid, since its unit $R^{\wedge}_I$ is not compact unless $I$ is nilpotent.
	\end{example}
	% CHECK: "in general not compactly generated" is not stated by Efimov. If it were compactly generated it would be Ind(Perf(R^wedge_I)) = D(R^wedge_I), since its compact objects are Perf(R^wedge_I); justify that this fails, e.g. for R = Z_p, I = (p).

	\begin{example}[Banach spaces as nuclear modules]
		\label[example]{ex:Nuclear_Banach}
		Heuristically, nuclear modules are ``Banach-type'' objects, as opposed to the compact objects such as $\prod_{\N}\Z_p$, which are of ``profinite'' or ``compact'' type. For instance, for every profinite set $S$ the module $C(S,\Z_p)$ of continuous functions is nuclear over $\Z_p^{\square}$ \cite[Lemma~3.8]{Andreychev_Thesis}. Taking $S = \N \cup \{\infty\}$, it follows that the $p$-adic Banach space $c_0(\N,\Z_p)$ of null sequences, a direct summand of $C(S,\Z_p)$, is nuclear, and inverting $p$ gives nuclear solid $\Q_p$-vector spaces such as $c_0(\N,\Q_p)$. In contrast, the compact object $\prod_{\N}\Z_p$ of $\D(\mathrm{Solid}_{\Z_p})$ is not nuclear: a compact nuclear object is dualizable, and the dualizable objects of $\D(\mathrm{Solid}_{\Z_p})$ are just the perfect complexes of $\Z_p$-modules \cite[Theorem~14.10]{Scholze_Analytic_Geometry}.
	\end{example}

	\begin{example}[Liquid and gaseous modules]
		\label[example]{ex:Nuclear_Liquid_Gaseous}
		Over the real and complex numbers, the solid theory is replaced by the \emph{liquid} theory: for $0 < p \leq 1$ there is a compactly generated symmetric monoidal category $\D(\mathrm{Liq}_p(\R))$ of $p$-liquid $\R$-vector spaces \cite{Scholze_Analytic_Geometry}, and one can form $\Nuc(\R^{p\text{-}\mathrm{liq}})$. The nuclear objects of $\D(\mathrm{Liq}_p)$ do not depend on $p$ \cite[Corollary~8.18]{Clausen_Scholze_Complex_Geometry}, and according to Clausen's lectures the rigidification of the category of liquid $\C$-vector spaces can be described in terms of operator algebras:
		\[
			\D^{\mathrm{liq}}(\C)^{\rig} \simeq \D\big(B(\Hh) \to B(\Hh)/S_{0+}\big),
		\]
		where $S_{0+} = \bigcap_{p > 0} S_p$ is the ideal of operators which are $p$-summable for all $p > 0$, cf.\ \Cref{ex:Homological_Epis}(3) and \Cref{ex:Compact_Operators_Intro}. The \emph{gaseous} theory of Clausen--Scholze \cite{Clausen_Scholze_Complex_Geometry} provides yet another analytic ring structure on $\C$, with a category of nuclear modules $\Nuc(\C^{\mathrm{gas}})$. In all cases the resulting categories are rigid, so dualizable but not compactly generated, and their continuous K-theory is the natural candidate for the ``analytic K-theory'' of $\R$ and $\C$, cf.\ \Cref{sec:First_Computations}.
	\end{example}
	% CHECK: the formula for D^liq(C)^rig is quoted from Clausen's Lecture 3 without proof.

	\begin{remark}
		\label[remark]{rem:Nuclear_Descent}
		The reason for working with $\Nuc$ rather than with all complete (solid) modules is that the former has much better formal properties: for instance, it is rigid, hence dualizable with a good continuous K-theory, and $\operatorname{Spa}(A,A^+) \mapsto \Nuc(A)$ satisfies descent for the analytic topology on affinoid adic spaces (Andreychev), see \Cref{sec:First_Computations}.
	\end{remark}
